% Generated draft: requires mathematical and source review.
\documentclass[11pt]{article}
\usepackage[margin=25mm]{geometry}
\usepackage{fontspec}
\setmainfont{Latin Modern Roman}
\usepackage{amsmath,amssymb,mathtools,unicode-math}
\setmathfont{Latin Modern Math}
\usepackage{enumitem,needspace,xurl,hyperref}
\setlength{\parindent}{0pt}
\setlength{\parskip}{5pt}
\setlength{\emergencystretch}{3em}
\newcommand{\N}{\mathbb{N}}
\newcommand{\Z}{\mathbb{Z}}
\newcommand{\Q}{\mathbb{Q}}
\newcommand{\R}{\mathbb{R}}
\newcommand{\C}{\mathbb{C}}
\newcommand{\F}{\mathbb{F}}
\newcommand{\A}{\mathbb{A}}
\newcommand{\PP}{\mathbb{P}}
\newcommand{\E}{\mathbb{E}}
\newcommand{\eps}{\varepsilon}
\newcommand{\dd}{\,\mathrm{d}}
\begin{document}
\section*{UnsolvedMath ID: 800280}
\textbf{Source problem code:} catalogue entry
\subsection{Definitions and mathematical statement}
Interpret a five-dimensional CR surface as an abstract CR manifold of hypersurface type and CR dimension two. Under the usual smooth interpretation, this consists of a smooth real manifold \(M\) of dimension \(5\) and a smooth complex subbundle \(T^{0,1}M\subset \C\otimes_{\R}TM\) of complex rank \(2\), satisfying
\[
T^{0,1}M\cap\overline{T^{0,1}M}=\{0\},\qquad
[\Gamma(T^{0,1}M),\Gamma(T^{0,1}M)]\subseteq\Gamma(T^{0,1}M).
\]
Write \(T^{1,0}M=\overline{T^{0,1}M}\), and let \(H\subset TM\) be the real rank-four distribution whose complexification is \(T^{1,0}M\oplus T^{0,1}M\). Strong pseudoconvexity means that near every point there is a nonvanishing real one-form \(\theta\) with \(\ker\theta=H\) for which the Levi form
\[
\mathcal L_{\theta}(Z,\overline W)=-i\,\dd\theta(Z,\overline W),\qquad Z,W\in T^{1,0}M,
\]
is positive definite, after changing the sign of \(\theta\) if necessary.

Under these conventions, the question is: for every such strongly pseudoconvex \(M\) and every \(p\in M\), does there exist an open neighborhood \(U\) of \(p\) and a smooth embedding
\[
F=(f_1,f_2,f_3):U\longrightarrow\C^3
\]
such that \(L f_j=0\) for every smooth local section \(L\) of \(T^{0,1}M|_U\) and every \(j\in\{1,2,3\}\)? Here an embedding is a homeomorphism onto its image with injective real differential; its image is consequently a real hypersurface in \(\C^3\).

The supplied statement does not specify differentiability, explicitly state hypersurface type, or name an ambient complex dimension. The formulation above makes the conventional smooth hypersurface interpretation explicit; these conventions require confirmation against the source. The term ``surface'' is retained from the source and is interpreted using its explicit real dimension \(5\).
\subsection{Short English statement}
Does every strongly pseudoconvex abstract CR manifold of real dimension \(5\) admit a local CR embedding near each point? Under the smooth hypersurface interpretation, the target is \(\C^3\). The source leaves the regularity and embedding conventions unspecified.
\subsection{Sources}
\begin{itemize}
\item Local supplied UnsolvedMath registry record, permanent ID \(800280\), exact source problem code \texttt{catalogue entry}; title: \emph{Several Complex Variables Problem List --- Problem 80}. The registry attributes the problem to Buzzard, Fornaess, and Krantz, with proposed year \(1995\), and labels it ``open.'' This is a supplied status label, not an independently verified mathematical or literature assessment.
\item Supplied source attribution: Buzzard, Fornaess, Krantz, \emph{The Several Complex Variables Problem List} (1995), Problem 80. Supplied URL: \url{https://arxiv.org/abs/math/9509202}. The record reports access on \texttt{2026-10-16} and extraction by \texttt{arxiv\_latex}; these are supplied provenance statements, not independently verified access or source checks.
\end{itemize}
\end{document}
